% 18 November 2026. Prepared for Edward Baker with OpenAI Codex LLM assistance.
% Actual model and effort unavailable; inherited helper settings. E26 source integration.
\section{A coupled hierarchy indexed by the adjacent pair}
\label{sec:lattice}

The split-pole moments can be attached to every adjacent pair of the
same divisor. This adds an index locating the pair, while imposing
compatibility between rows that share a zero. The resulting description
retains correlations which independent moment rows would discard.
We develop the formulation proposed in \cite{BakerCM29}, using the
single-pair identities of \cite{BakerCM08,BakerCM09}. The conclusions
are exact representation and error identities. An estimate preventing
concentration into an individual gap remains a separate requirement.

\subsection{Edge coordinates and the complete gap interaction}

First suppose the complete divisor is real and simple on the time
interval under consideration, and order its zeros as $x_n(t)$,
$n\in\mathbb Z$. The integer labels are continued locally in time.
All infinite-lattice assertions below are componentwise; uniform
estimates concern a fixed finite set of edges on a compact interval
before the relevant loss of simplicity or separation. Define
\begin{equation}
 g_n=x_{n+1}-x_n>0,\qquad
 c_n=\frac{x_n+x_{n+1}}2,\qquad D_n=\frac{g_n^2}{4},
 \label{lat:coordinates}
\end{equation}
and, in the orientation of \eqref{split:def},
\begin{equation}
 M_n(p,q;t)=\sum_{k\ne n,n+1}
       (x_k-x_{n+1})^{-p}(x_k-x_n)^{-q},
 \quad p,q\ge0,\quad p+q\ge2.
 \label{lat:moments}
\end{equation}
The index $n$ belongs to the edge $(x_n,x_{n+1})$. It does not
specify a new exterior source. In particular,
$M_n(1,1)$ is the interaction $M$ of that edge,
$M_n(r,r)=S_r^{(n)}$, and
$\bigl(M_n(r-1,r)+M_n(r,r-1)\bigr)/2=T_r^{(n)}$ for $r\ge2$.

For these exact identities, a strictly separated count-two contour
suffices at each edge; the stronger collar ratio in
\eqref{eq:domain} is not imposed on a general adjacent-edge lattice.
On a compact subinterval, a finite cover by such contours gives
positive denominator bounds for each retained edge. The finite
algebra in Proposition~\ref{split:prop} and its complete-product
passage therefore apply with those bounds. Its recurrence
\eqref{split:mixed}--\eqref{split:axes} holds at every $n$, replacing
$P_{p,q}$ by $M_n(p,q)$ and $d_0$ by $g_n$. In particular, in
forward physical time,
\begin{equation}
 g_n'=\frac4{g_n}-2g_nM_n(1,1),\qquad
 D_n'=2-4D_nM_n(1,1).
 \label{lat:gap-law}
\end{equation}
There is no additional nearest-neighbor term to append to the
single-pair recurrence. Coupling in $n$ is already encoded by the
common zeros and the identities derived below.

The complete sums in \eqref{lat:moments} are absolute at their stated
orders. For targets in a fixed compact spatial set and $N$ sufficiently
large, the terms with $|\rho|>N$ are bounded by
$C_{p,q}|\rho|^{-p-q}$. The unit-count bound
\eqref{eq:count-safe} consequently gives the tail
\begin{equation}
 \sum_{|\rho|>N}
 |\rho-x_{n+1}|^{-p}|\rho-x_n|^{-q}
 =O_{p,q,I,\mathcal B}(N^{1-p-q}\log N),
 \label{lat:tail}
\end{equation}
uniformly for a fixed finite edge block $\mathcal B$ with the stated
clearances. Opposite and conjugate occurrences and all multiplicities
are included. No uniformity in a growing block, height or moment order
is inferred from this estimate.

\begin{proposition}[All-gap expression]\label{lat:all-gap-prop}
When the complete divisor is real and ordered, put
\begin{equation}
 L_{n,m}=\sum_{j=1}^{m}g_{n-j},\qquad
 R_{n,m}=\sum_{j=1}^{m}g_{n+j},\qquad m\ge1.
 \label{lat:distances}
\end{equation}
Then the complete interaction is
\begin{equation}
 M_n(1,1)=\sum_{m\ge1}
 \left\{\frac1{L_{n,m}(L_{n,m}+g_n)}
       +\frac1{R_{n,m}(R_{n,m}+g_n)}\right\}.
 \label{lat:all-gap}
\end{equation}
Thus \eqref{lat:gap-law} and \eqref{lat:all-gap} describe a nonlocal
evolution of the complete gap sequence. They do not specify the
absolute translation coordinate.
\end{proposition}

\begin{proof}
The left exterior zero $x_{n-m}$ has distances $-L_{n,m}$ and
$-(L_{n,m}+g_n)$ from the left and right targets. The right exterior
zero $x_{n+1+m}$ has distances $R_{n,m}+g_n$ and $R_{n,m}$.
These exhaust the exterior divisor. Substitution in
\eqref{lat:moments} proves the formula; absolute convergence from
\eqref{lat:tail} permits the partition. The gap law has already
eliminated the common center velocity, so it determines relative
motion only.
\end{proof}

Every gap enters many rows through cumulative distances in
\eqref{lat:distances}. This is a spatially nonlocal lattice, rather
than a closed nearest-neighbor equation. For the stationary cosine
control with every gap equal to $g$, each side of
\eqref{lat:all-gap} is
$g^{-2}\sum_{m\ge1}[m(m+1)]^{-1}=g^{-2}$. Hence
$M_n(1,1)=2/g^2$ and $g_n'=0$, checking the physical normalization.
This noncanonical control does not model the canonical population.

If only a finite consecutive real block is known, use
\eqref{lat:moments} with the remaining complete divisor, which may
be nonreal. The part of \eqref{lat:all-gap} belonging to known real
zeros is still valid, but the omitted contribution must be retained
as a complete exterior interaction. Nonreal groups cannot be replaced
by the positive real-gap terms of \eqref{lat:all-gap}. The row
identities and the compatibility statements below remain valid on
cleared charts of the real block in this setting.

\subsection{Overlap identities and a common meromorphic field}

\begin{proposition}[Adjacent-row compatibility]\label{lat:compat-prop}
For every integer $r\ge2$,
\begin{equation}
 M_n(r,0)+(-g_n)^{-r}
 =M_{n+1}(0,r)+g_{n+1}^{-r}.
 \label{lat:overlap}
\end{equation}
Within a row, for $p,q\ge1$ and $p+q\ge3$,
\begin{equation}
 g_nM_n(p,q)=M_n(p,q-1)-M_n(p-1,q).
 \label{lat:row-algebra}
\end{equation}
These are necessary constraints on rows arising from a common divisor.
They are not asserted to characterize all realizable moment arrays.
\end{proposition}

\begin{proof}
Both sides of \eqref{lat:overlap} equal
$\sum_{k\ne n+1}(x_k-x_{n+1})^{-r}$, with any nonreal exterior
occurrences also included. The left side restores $x_n$ to the
sum for edge $n$, and the right side restores $x_{n+2}$ to the sum
for edge $n+1$. All terms are absolute for $r\ge2$.
For \eqref{lat:row-algebra}, set
$X=\rho-x_{n+1}$, $Y=\rho-x_n=X+g_n$. The identity
\[
 X^{-p}Y^{-(q-1)}-X^{-(p-1)}Y^{-q}
 =(Y-X)X^{-p}Y^{-q}
\]
can be summed absolutely in the stated range. At $(p,q)=(1,1)$,
the same identity holds only as the common-cutoff difference of
the first-power sums: its combined summand is $g_n/(XY)$ and
decays quadratically. Separate absolute first-power sums are neither
required nor introduced.
\end{proof}

The common generating field is $B(t,z)=H_z(t,z)/H(t,z)$.
Removing the two target poles of edge $n$ gives
\begin{equation}
 R_n(t,z)=B(t,z)-\frac1{z-x_n(t)}-\frac1{z-x_{n+1}(t)}.
 \label{lat:resolvent}
\end{equation}
The removed singularities are removable at these simple zeros.
For three consecutive distinct labels, direct subtraction gives
\begin{equation}
 R_{n+1}(t,z)=R_n(t,z)+\frac1{z-x_n(t)}
                         -\frac1{z-x_{n+2}(t)}.
 \label{lat:resolvent-shift}
\end{equation}
This identity requires no reality assumption on the other zeros.
The tail conventions are those of the same complete field $B$ on
both sides; there is no independently chosen additive regularizer.

The Taylor coefficients of $R_n$ at the two targets recover the
pure-axis moments. For $r\ge2$,
\begin{equation}
 \frac{\partial_z^{r-1}R_n(t,x_{n+1})}{(r-1)!}
       =-M_n(r,0),\qquad
 \frac{\partial_z^{r-1}R_n(t,x_n)}{(r-1)!}
       =-M_n(0,r).
 \label{lat:axis-taylor}
\end{equation}
Indeed, differentiation of an exterior simple pole gives
$(-1)^{r-1}(z-\rho)^{-r}$; at the target this is
$-(\rho-z)^{-r}$. The derivatives remove the genus-one constant
regularizers, and the remaining series is absolute.

\begin{proposition}[Mixed generating function]\label{lat:generator-prop}
Write $a=x_n$, $b=x_{n+1}$ at a fixed time. In sufficiently small
cleared neighborhoods of the two targets,
\begin{equation}
 -\frac{R_n(b+u)-R_n(a+v)}{b-a+u-v}
 =\sum_{p,q\ge1}M_n(p,q)u^{p-1}v^{q-1}.
 \label{lat:generator}
\end{equation}
The quotient has its divided-difference interpretation when its
two arguments coincide within the analytic domain of $R_n$.
Moreover, off the zeros,
\begin{equation}
 B_t=-B_{zz}-2BB_z.
 \label{lat:burgers}
\end{equation}
\end{proposition}

\begin{proof}
For one exterior pole,
\[
 -\frac{(b+u-\rho)^{-1}-(a+v-\rho)^{-1}}
        {b-a+u-v}
 =\frac1{(\rho-b-u)(\rho-a-v)}.
\]
The constant regularizers cancel in the difference. Choose the
neighborhoods so that $|u|/|\rho-b|$ and $|v|/|\rho-a|$ are
uniformly bounded by a number less than one for every exterior
occurrence. Expanding the two geometric series gives
\eqref{lat:generator}. Their absolute sum is bounded by a constant
times $\sum_\rho|\rho-b|^{-1}|\rho-a|^{-1}$, which is finite by
the fixed-order tail estimate and local clearance. This justifies
the interchange of the two power series and the complete divisor.
If the arguments coincide, the left side is $-R_n'$ by analytic
continuation of the divided difference. Finally,
$H_t/H=-(B_z+B^2)$ follows from the heat equation. Differentiate
in $z$ to obtain \eqref{lat:burgers}.
\end{proof}

Equations \eqref{lat:resolvent-shift}--\eqref{lat:burgers} express
all rows through one field. A Taylor translation cannot cross an
unremoved neighboring pole: use overlapping cleared neighborhoods
or remove the entire chosen block first. The meromorphic equation
is another exact expression of the original heat evolution, not an
additional estimate on its canonical solution. An arbitrary array
of individually evolved or truncated moment rows need not satisfy
these overlap identities or arise from any divisor.

\subsection{Node forces, block cancellation and its error budget}

At a simple root define the regular part
\begin{equation}
 U_n(t)=\lim_{z\to x_n}
       \left(B(t,z)-\frac1{z-x_n}\right)
       =\frac{H_{zz}(t,x_n)}{2H_z(t,x_n)}.
 \label{lat:node-force}
\end{equation}
Taylor expansion of $H$ at $x_n$ gives the second equality. Implicit
differentiation of $H(t,x_n(t))=0$ and $H_t=-H_{zz}$ gives
$x_n'=2U_n$. This definition uses a finite part of the complete
source, not an unregularized infinite first-power sum.

\begin{proposition}[Block flux identity]\label{lat:flux-prop}
For a fixed finite block of edges $n=a,\ldots,b$ on a common
simple-root interval,
\begin{align}
 g_nM_n(1,1)&=\frac2{g_n}-(U_{n+1}-U_n),
                         \label{lat:edge-flux}\\
 \sum_{n=a}^{b}g_nM_n(1,1)
 &=2\sum_{n=a}^{b}\frac1{g_n}-(U_{b+1}-U_a).
                         \label{lat:block-flux}
\end{align}
The same formulas hold for a finite real block with nonreal exterior,
provided the complete interactions and source finite parts are used.
\end{proposition}

\begin{proof}
The node law gives $g_n'=2(U_{n+1}-U_n)$. Comparison with
\eqref{lat:gap-law} proves \eqref{lat:edge-flux}. Summing the
successive differences telescopes exactly, giving
\eqref{lat:block-flux}. No positivity of exterior summands enters
these two operations.
\end{proof}

The cancellation has a precise, limited consequence for error
propagation. Put $A_{a,b}=\sum_{n=a}^b g_nM_n(1,1)$. If the
gaps are exact and common approximations to the endpoint forces
satisfy $|\widehat U_j-U_j|\le\varepsilon_j$ for
$j=a,b+1$, reconstruct the aggregate by the right side of
\eqref{lat:block-flux}. Then
\begin{equation}
 |\widehat A_{a,b}-A_{a,b}|
 \le\varepsilon_a+\varepsilon_{b+1}.
 \label{lat:exact-gap-error}
\end{equation}
For a uniform node error $\varepsilon$, this is $2\varepsilon$,
independent of the number of edges. Bounding each edge-force
difference separately and summing absolute allowances would pay
up to $2(b-a+1)\varepsilon$. The smaller aggregate allowance uses
the same approximated node at each shared endpoint, so its interior
errors cancel before absolute values are taken.

If the gaps are also approximated, assume $\widehat g_n>0$ and
$g_n>0$, and define
$\widehat A_{a,b}=2\sum_{n=a}^b\widehat g_n^{-1}
-(\widehat U_{b+1}-\widehat U_a)$. The full bound is
\begin{align}
 |\widehat A_{a,b}-A_{a,b}|
 &\le\varepsilon_a+\varepsilon_{b+1}
       +2\sum_{n=a}^b
          \frac{|\widehat g_n-g_n|}{\widehat g_n g_n}
                                                        \notag\\
 &\le\varepsilon_a+\varepsilon_{b+1}
       +2\sum_{n=a}^b
          \frac{|\widehat g_n-g_n|}
               {\min(\widehat g_n,g_n)^2}.
 \label{lat:uncertain-gap-error}
\end{align}
This follows from the exact difference of two reciprocals. Certified
positive lower gap bounds may replace the denominators in the final
allowance. Small gaps can make this payment large. These bounds
concern an aggregate reconstructed by the identity; independently
approximating each product $g_nM_n(1,1)$ requires its additional
weight and moment errors.

The same-information direct comparator has exactly the same
cancellation. The block length satisfies
\begin{equation}
 \frac{\mathrm d}{\mathrm dt}(x_{b+1}-x_a)
       =2(U_{b+1}-U_a)
       =4\sum_{n=a}^b\frac1{g_n}-2A_{a,b}.
 \label{lat:block-length}
\end{equation}
Thus \eqref{lat:exact-gap-error} is a correlated-error advantage
over an independently summed allowance, not a demonstrated gain
over direct common-source evolution. A block can preserve its total
length while one interior gap decreases to zero. Neither
\eqref{lat:block-flux} nor its endpoint-error bound protects the
regional minimum. Such protection would need a further estimate
that prevents concentration of the defect into one edge.

\subsection{Compatible reconstruction and the collision boundary}

One finite polynomial heat field produces all its roots and moment
rows at once. Before the relevant collisions, its rows satisfy the
gap, overlap, resolvent and block identities above exactly. This is
the compatibility supplied by the shared finite-divisor construction
of \cite{BakerCM15}, developed further in Section~\ref{sec:finite}.
Its higher moments are those of the same finite divisor; they are
not independently set to zero. Independent truncation in $(p,q)$
at each edge is therefore a different construction, generally
without a common underlying divisor.

Compatibility alone does not establish approximation to the
canonical field. Sections~\ref{sec:common-time}--\ref{sec:counted-transfer}
prove a common-time finite-product theorem with one initialization and
explicit field, jet and fixed-moment errors. A finite-block application
must supply a compact set containing every relevant contour, a common
field error, matching labels and the minimum clearance across that block.
The pairwise transfer then applies with all these constants paid,
provided every retained pair satisfies the chart and strong-collar
hypotheses of \eqref{eq:domain} and Section~\ref{sec:counted-transfer}. No uniformity
as the block size, height or moment order grows is included. An
index cutoff also leaves complete nonlocal exterior forces and
boundary data; it cannot be implemented by dropping the omitted
zeros without an error estimate.

The adjacent-edge charts have a further limitation at mergers.
If $x_n=x_{n+1}$, the symmetric $S/T$ coordinates of that target
pair can remain regular, as in Section~\ref{sec:integer}. For the
neighboring edge $(x_{n-1},x_n)$, however, the exterior occurrence
$x_{n+1}$ meets a target; the analogous failure occurs for
$(x_{n+1},x_{n+2})$. Their exterior separation is lost. It is
therefore incorrect to infer that every adjacent split-pole chart
extends regularly through the same merger. A shared polynomial,
or suitable larger-cluster symmetric coefficients, can remain
regular while these individual charts must stop or be replaced.

Physical quadratic and cosine-plus-constant collision controls
satisfy the applicable simple-root identities before collision.
The lattice representation does not exclude those events. Its
role is to preserve the common-divisor constraints and expose
correlations such as \eqref{lat:block-flux}; it supplies neither
a universal positivity principle nor stability of a finite
closure. The next quantitative question is a common-source
estimate controlling concentration into individual gaps over a
specified physical interval, with complete exterior and boundary
payments and the direct evolution given the same information.
